A parametric surface is described by a function $\bm{r}(u,v): \mathbb{R}^2 \rightarrow \mathbb{R}^3 = x(u,v)\bm{\hat{\imath}} + y(u,v)\bm{\hat{\jmath}} + z(u,v)\bm{\hat{k}}$, defined over $u \in [u_0, u_1]$ and $v \in [v_0, v_1]$.
It is assumed that $\bm{r}(u,v)$ is of class $C^2$, so the partial derivatives $\displaystyle \bm{r}_u = \frac{\partial r}{\partial u}$, $\displaystyle \bm{r}_v = \frac{\partial r}{\partial v}$,
$\bm{H}_x = \nabla \otimes \nabla x$, $\bm{H}_y = \nabla \otimes \nabla y$, and $\bm{H}_z = \nabla \otimes \nabla z$ exist. They can be provided by the user or calculated numerically.

Because of the extensive use of derivative and integral operators, repeatedly calculating some properties of similar shapes can be expensive. All parametric surfaces have shared pointer to a common base shape
$\bm{r}(u,v)$, provided that they are defined by a set related parametric equations $\bm{r}'(u,v) = \bm{M}\bm{r}(u,v) + \bm{r}_0$, where $\bm{M}$ is an invertible linear transformation matrix in $\mathbb{R}^{3\times3}$,
and $\bm{r}_0$ is a translation vector.

\subsection{Extreme points}
Determining the extreme points in any Cartesian direction is a constrained optimization problem, since they must occur at either at the boundary or a critical point within. Take the x-coordinate of a base parametric
surface for example, this involves locating the global minimum of $x(u,v)$ on the box interval $[u_0, u_1] \times [v_0, v_1]$.

Because of the low dimensionality of the functions (2), computing the Hessian (of $x(u,v)$, for example) is cheap. The method of choice is the projected Newton's method. Similar to the Levenberg–Marquardt least-square
algorithm, an addition of $\alpha \bm{I}$ (for an $\alpha \geq 0$) to the Hessian may be needed to ensure that it is positive definite. At every iteration, a backtracking line search along the search direction
$(\bm{H} + \alpha\bm{I})^{-1} \nabla x$ is needed to ensure that the solution stays constrained while satisfying the Armijo condition.

For a transformed shape, the x-component and its derivative functions are given as
\begin{align}
    x'(u,v) &= \bm{M}_0\bm{r}(u,v) + x_0 \\
    \left(\nabla x'\right)^T &= \bm{M}_0
    \begin{bmatrix}
        \bm{r}_u & \bm{r}_v
    \end{bmatrix}
    \\
    \bm{H}_{x'} &= M_{00}\bm{H}_x + M_{01}\bm{H}_y + M_{02}\bm{H}_z
\end{align}
where $\bm{M}_0$ is the row vector corresponding to the first row of $\bm{M}$.

If $\bm{M}$ is a scaling matrix that scales the x-coordinate by a scalar $k_x$, then the extrema points of $x'$ are $x'_{\text{min}} = k_x x_{\text{min}} + x_0$ and $x'_{\text{max}} = k_x x_{\text{max}} + x_0$.

Similar procedures can be followed for the y and z coordinates.

\subsection{Surface area}
The surface area of a base parametric surface is
\begin{equation}
    S = \int_{v_0}^{v_1} \int_{u_0}^{u_1} \Vert \bm{r}_u \times \bm{r}_v \Vert du dv.
    \label{eqn:base_parametric_surface_area}
\end{equation}

Rotation and reflection do not alter the surface area. Uniform scaling by a scalar $K$ scales the surface area by $K^2$. Any other transformation will require a re-evaluation of the integral in Equation 
\ref{eqn:base_parametric_surface_area} as
\begin{equation}
    S = \int_{v_0}^{v_1} \int_{u_0}^{u_1} \Vert \bm{M}\bm{r}_u \times \bm{M}\bm{r}_v \Vert du dv.
\end{equation}

\subsection{Volume}
The volume is only defined for a closed surface, since we can take advantage of the Divergence Theorem
\begin{equation}
    \iiint_\Gamma \left( \nabla \cdot \bm{F} \right) dV = \oiint_{\partial\Gamma} \bm{F} \cdot d\bm{S}.
\end{equation}

Rather than integrating over a volume, the better alternative is to find an $\bm{F}$ such that $\nabla \cdot \bm{F} = 1$ and integrate it over the surface. From Equation \ref{eqn:base_parametric_surface_area},
if we follow the convention
\begin{equation}
    S = \oiint_{\partial \Gamma} dS = \oiint_{\partial \Gamma} \bm{\hat{n}} \cdot d\bm{S},
\end{equation}
the unit normal vector $\bm{\hat{n}}$ and the differential surface area $d\bm{S}$ can be defined as
\begin{align}
    \bm{\hat{n}} &\equiv \frac{\bm{r}_u \times \bm{r}_v}{\Vert \bm{r}_u \times \bm{r}_v \Vert} \label{eqn:n_parametric} \\
    d\bm{S} &\equiv \left( \bm{r}_u \times \bm{r}_v \right) dudv. \label{eqn:dS_parametric}
\end{align}

The last task to to find a function $\bm{F}$ such that $\nabla \cdot \bm{F} = 1$. The most obvious candidates are $\bm{F} = x \bm{\hat{\imath}}$, $\bm{F} = y \bm{\hat{\jmath}}$, and $\bm{F} = z \bm{\hat{k}}$.
Choosing the last candidate, the expression for the volume is:
\begin{align*}
    V &= \int_{v_0}^{v_1} \int_{u_0}^{u_1} z \hat{k} \cdot \left( \bm{r}_u \times \bm{r}_v \right) dudv \\
    &= \int_{v_0}^{v_1} \int_{u_0}^{u_1} z \left( \frac{\partial x}{\partial u}\frac{\partial y}{\partial v} - \frac{\partial x}{\partial v}\frac{\partial y}{\partial u} \right) dudv. \numberthis
    \label{eqn:V_parametric}
\end{align*}

The volume of a transformed shape is that of the base shape scaled by $\lvert \det \bm{M} \rvert$.

\subsection{Checking if a point on the surface}
A point $\bm{p}$ is on the parametric surface $\bm{r}(u,v)$ (regardless of if it's a base or transformed surface) if there exists a pair $(u,v)$ that satisfies the system of equation
\begin{equation}
    \bm{p} - \bm{r}(u,v) = \bm{0}.
    \label{eqn:point_on_parametric_surface}
\end{equation}

This is a non-linear least square problem, since the system has 3 equations but only 2 unknowns. Given the least-square solution $(u^*, v^*)$ to Equation \ref{eqn:point_on_parametric_surface}, then $\bm{p}$ is on
the parametric surface if
\begin{equation}
    \Vert \bm{p} - \bm{r}(u^*,v^*) \Vert = 0.
\end{equation}
In other words, $(u^*, v^*)$ has to satisfy Equation \ref{eqn:point_on_parametric_surface} exactly (or within a numerical tolerance) rather than merely minimizing the residual.

\subsection{Checking if a point is inside}
The concept of "inside" (and by extension "outside") is only well-defined if the surface is closed. For simplicity, the surface is assumed convex. A point $\bm{p}$ is inside the volume enclosed by the parametric
surface if regardless of the direction at which it travels, there is a positive distance to the surface. This problem reduces to two distance calculation problems, one with a direction $\bm{\hat{\omega}}$ and the
other with $-\bm{\hat{\omega}}$, and checking whether both distances are positive. The choice of $\bm{\hat{\omega}}$ is arbitrary. Refer to section \ref{sect:distance_to_parametric_surface} for the actual distance
calculation.

\subsection{Checking if another volume overlaps}
\subsection{Distance to surface}
\label{sect:distance_to_parametric_surface}
Given a point $\bm{p}$ travelling at direction $\bm{\hat{\omega}}$, the distance $s$ to surface of the sphere satisfies the equation
$$
    \bm{p} + \bm{\hat{\omega}}s = \bm{M}\bm{r}(u,v) + \bm{r}_0.
$$

This can be arranged into an objective function of three variables:
\begin{equation}
    \bm{F}(u,v,s) = \bm{p}_0 + \bm{\hat{\omega}}s - \bm{M}\bm{r}(u,v) - \bm{r}_0,
\end{equation}
and the goal is to solve for the solution vector $\bm{s} = [u, v, s]^T$ such that $\bm{F}(\bm{s}) = \bm{0}$.

Since the Jacobian of $\bm{F}$ is known
\begin{equation}
    \bm{J} =
    \begin{bmatrix}
        -\bm{M}\bm{r}_u & -\bm{M}\bm{r}_v & \bm{\hat{\omega}}
    \end{bmatrix},
\end{equation}
Newton's method can be used to solve for $\bm{s}$. The solution is permissible if $u$ and $v$ are both in bound, and $s \geq 0$.

\subsection{Outward normal vector}
The outward normal vector is defined in Equation \ref{eqn:n_parametric}.

\section{Shape defined by an implicit function}